
\medskip

\emph{Principaux domaines abordés}: Probabilités conditionnelles et indépendance. Variables aléatoires.

\bigskip

\begin{center}\emph{Les deux parties de cet exercice sont indépendantes.}\end{center}

\textbf{Partie 1}

\medskip

%Julien doit prendre l'avion; il a prévu de prendre le bus pour se rendre à l'aéroport. 
%
%S'il prend le bus de 8 h, il est sûr d'être à l'aéroport à temps pour son vol.
%
%Par contre, le bus suivant ne lui permettrait pas d'arriver à temps à l'aéroport.
%
%Julien est parti en retard de son appartement et la probabilité qu'il manque son bus est de $0,8$.
%
%S'il manque son bus, il se rend à l'aéroport en prenant une compagnie de voitures privées; il a alors une probabilité de $0,5$ d'être à l'heure à l'aéroport.
%
%\smallskip
%
%On notera :
%
%\setlength\parindent{12mm}
%\begin{itemize}
%\item[$\bullet~~$] $B$ l'évènement: \og Julien réussit à prendre son bus \fg ;
%\item[$\bullet~~$] $V$ l'évènement: \og Julien est à l'heure à l'aéroport pour son vol \fg.
%\end{itemize}
%\setlength\parindent{0mm}
%
%\medskip

\begin{enumerate}
\item %Donner la valeur de $P_B(V)$.
S'il a pris le bus il est à l'heure pour son vol, donc $P_B(V) = 1$.

\item On représente la situation par un arbre pondéré.

\begin{center}
\pstree[treemode=R,nodesepB=3pt,levelsep=3cm]{\TR{}}
{\pstree{\TR{$B$~~} \taput{$0,2$}}
	{\TR{$V$}\taput{1}
	\TR{$\overline{V}$}\tbput{0}
	}
\pstree{\TR{$\overline{B}$~~} \tbput{$0,8$}}
	{\TR{$V$}\taput{0,5}
	\TR{$\overline{V}$}\tbput{0,5}
	}	
}
\end{center}

\item %Montrer que $P(V) = 0,6$.
D'après la loi des probabilités totales :

$P(V) = P(B \cap V) + P\left(\overline{B} \cap V\right) = 0,2 \times 1 + 0,8 \times 0,5 = 0,2 + 0,4 = 0,6$.

\item %Si Julien est à l'heure à l'aéroport pour son vol, quelle est la probabilité qu'il soit arrivé à l'aéroport en bus ? Justifier.
On calcule $P_V(B) = \dfrac{P(V \cap B)}{P(V)} = \dfrac{P(B \cap V)}{P(V)} = \dfrac{0,2}{0,6} = \dfrac13$.
\end{enumerate}

\bigskip

\textbf{Partie 2}

\medskip

%Les compagnies aériennes vendent plus de billets qu'il n'y a de places dans les avions car certains passagers ne se présentent pas à l'embarquement du vol sur lequel ils ont réservé.
%On appelle cette pratique le surbooking.
%
%Au vu des statistiques des vols précédents, la compagnie aérienne estime que chaque passager a 5\,\% de chance de ne pas se présenter à l'embarquement.
%
%Considérons un vol dans un avion de $200$~places pour lequel $206$~billets ont été vendus.
%On suppose que la présence à l'embarquement de chaque passager est indépendante des autres passagers et on appelle $X$ la variable aléatoire qui compte le nombre de passagers se présentant à l'embarquement.
%
%\medskip

\begin{enumerate}
\item %Justifier que $X$ suit une loi binomiale dont on précisera les paramètres.
La présence à l'embarquement de chaque passager est indépendante des autres et chaque passager a la même probabilité 0,95 d'être présent, donc la variable aléatoire $X$ suit la loi binomiale de paramètres $n = 206$ et $p = 0,95$.
\item %En moyenne, combien de passagers vont-ils se présenter à l'embarquement ?
On a E $ = n \times p = 206 \times 0,95 = 195,7 \approx 196$.

En moyenne sur 206 titulaires d'un billet à peu près 196 vont se présenter.
\item %Calculer la probabilité que $201$ passagers se présentent à l'embarquement. Le résultat sera arrondi à $10^{-3}$ près.
On a $P(X = 201) = \binom{206}{201} \times 0,95^{201} \times 0,05^{5} \approx \np{0,03063}$, soit 0,031 au millième près.
\item %Calculer $P(X \leqslant 200)$, le résultat sera arrondi à $10^{-3}$ près. Interpréter ce résultat dans le contexte de l'exercice.
La calculatrice donne $P(X \leqslant 200) \approx \np{0,9477}$, soit 0,948 au millième près.

 Il est donc à peu près certain que l'avion sera au mieux juste rempli.
\item %La compagnie aérienne vend chaque billet à $250$ euros.

%Si plus de $200$ passagers se présentent à l'embarquement, la compagnie doit rembourser le billet d'avion et payer une pénalité de $600$ euros à chaque passager lésé.

%On appelle :

%\setlength\parindent{12mm}
%\begin{description}
%\item[ ] $Y$ la variable aléatoire égale au nombre de passagers qui ne peuvent pas embarquer
%bien qu'ayant acheté un billet;
%\item[ ] $C$ la variable aléatoire qui totalise le chiffre d'affaire de la compagnie aérienne sur ce vol.
%\end{description}

On admet que $Y$ suit la loi de probabilité donnée par le tableau suivant:

\begin{center}
\begin{tabularx}{\linewidth}{|c|*{7}{>{\centering \arraybackslash}X|}}\hline
$y_i$					&0	&1	&2	&3	&4	&5	&6\\ \hline
$P\left(Y = y_i\right)$	&\np{0,94775}&\np{0,03063}&\np{0,01441}&\np{0,00539}&\np{0,00151}&\np{0,00028}&$\blue\np{0,00003}$\\ \hline
$c_i$&$\blue\np{51500}$&$\blue\np{50650}$&$\blue\np{49800}$&$\blue\np{48950}$&$\blue\np{48100}$&$\blue\np{47250}$&$\blue\np{46400}$\\ \hline
\end{tabularx}
\end{center}

	\begin{enumerate}
		\item %Compléter la loi de probabilité donnée ci-dessus en calculant $P(Y = 6)$.
En complétant à 1 la somme des probabilités données dans le tableau on trouve:\\
 $P(Y = 6) =1 - \left ( P(Y=0) + P(Y=1) + \cdots + P(Y=5)\strut\right )= \np{0,00003}$.
 
		\item %Justifier que: $C = \np{51500} - 850Y$.
La compagnie a encaissé $206 \times 250 = \np{51500}$~\euro{} et elle devra rembourser 850~\euro{} à chaque client ne pouvant embarquer, donc $C = \np{51500} - 850Y$
		\item %Donner la loi de probabilité de la variable aléatoire $C$ sous forme d'un tableau.
Voir le tableau ci-dessus.

%Calculer l'espérance de la variable aléatoire $C$ à l'euro près.
On a E$(C) = \np{51500} \times \np{0,94775} + \ldots \np{46400} \times \np{0,00003} \approx \np{51429,2}$, soit \np{51429}~\euro{} à l'euro près

		\item %Comparer le chiffre d'affaires obtenu en vendant exactement $200$ billets et le chiffre d'affaires moyen obtenu en pratiquant le surbooking.
En vendant exactement 200 billets, la compagnie fera un chiffre d'affaires de $200\times 250$ soit $\np{50000}$ euros.

En pratiquant le surbooking, la compagnie peut espérer un chiffre d'affaires de \np{51429} euros.
	\end{enumerate}
\end{enumerate}

\bigskip

